\section{Related Work}
\label{sec:related_work}

\subsection{Intrusion Detection Systems}
Traditional IDSs rely on signature matching, which remains effective for known attack templates but fails on zero-day or rapidly evolving threats. Recent surveys show that learning-based IDSs now span classical machine learning, deep neural networks, graph-based methods, and large-scale IoT-oriented pipelines \cite{liao2024survey, moustafa2023explainable}. At the same time, dataset quality remains a major bottleneck: recent reviews emphasize that many public NIDS benchmarks contain outdated traffic patterns, shortcut features, or split assumptions that can overestimate deployment performance \cite{datasetreview2025}. This issue is directly relevant to our work because the empirical conclusions of Bi-ARL change noticeably across NSL-KDD, UNSW-NB15, and CIC-IDS2017. Importantly, these benchmarks have not disappeared from the literature: even recent transformer-based cross-dataset studies in 2025 still evaluate on NSL-KDD, UNSW-NB15, and CIC-IDS2017 while adding stronger calibration-oriented analysis \cite{xin2025crossdatasettransformer}.

\subsection{Generalization and Modern Deep IDS Models}
An important recent trend is to move beyond single-dataset closed-set evaluation. Golchin \emph{et al.} proposed SSCL-IDS, a self-supervised contrastive approach designed to improve detection under unseen traffic distributions \cite{golchin2024ssclids}. Xu \emph{et al.} further extended self-supervised learning to graph neural network based flow modeling, showing that representation learning can reduce the dependence on expensive labels while remaining competitive with supervised approaches \cite{xu2024gnnssl}. On the architecture side, transformer-based IDS models have also shown strong in-distribution performance; for example, IDS-MTran uses multi-scale branches and transformer backbones to obtain high accuracy on NSL-KDD, UNSW-NB15, and CIC-DDoS2019 \cite{xi2024idsmtran}. Xin and Xu further combine transformer-based IDS, calibration, and AUC-oriented optimization in a 2025 cross-dataset study, reinforcing that strong modern baselines should now be judged not only by headline accuracy but also by reliability under dataset shift \cite{xin2025crossdatasettransformer}. These works define a stronger modern comparison set than legacy RL-only baselines, but they still generally optimize static prediction models rather than explicitly incorporating an adaptive attacker into training.

\subsection{Robustness, Open-Set Detection, and Cross-Dataset Evaluation}
Robustness has emerged as a separate research axis from plain accuracy. Wang \emph{et al.} argue that adversarial perturbations and distribution shift should be treated as first-class evaluation criteria for ML-based NIDSs, rather than as optional stress tests \cite{wang2023robustnessnids}. Cantone \emph{et al.} provide direct evidence that models with strong within-dataset performance can deteriorate severely under cross-dataset testing, highlighting the gap between benchmark success and deployment reliability \cite{cantone2024crossdataset}. In parallel, VAEMax tackles open-set intrusion detection by combining OpenMax with a variational autoencoder to identify previously unseen attacks on CIC-IDS2017 and CSE-CIC-IDS2018 \cite{qiu2024vaemax}. Newer dataset efforts such as SHIELD/NF-UQ-NIDS and 5G-NIDD further indicate that the field is moving toward more contemporary, standardized, and domain-specific flow collections rather than relying only on older enterprise-network benchmarks \cite{shield2025, siriwardhana2025nidd}. These results motivate our focus on adaptive attacker-defender interaction, but they also imply that robustness claims should be made conservatively and backed by explicit protocol-sensitive evaluation.

\subsection{Adversarial Reinforcement Learning}
Reinforcement learning applied to security games allows for adaptive defense strategies and naturally accommodates sequential attacker-defender interaction. However, many practical RL-based IDS formulations still update the defender against either a static disturbance model or a simultaneously trained opponent, which can blur the distinction between ordinary multi-agent co-adaptation and a genuine best-response adversary. Our work focuses on this gap: rather than only using adversarial interaction as data augmentation, we formulate training as a nested optimization process in which the defender is optimized against an attacker that is repeatedly adapted toward a stronger response. This places the method closer to a Stackelberg-style robust control perspective than to standard simultaneous-update MARL. Empirically, however, our results also show that this formulation alone does not guarantee competitive performance on modern datasets, which is an important difference from stronger supervised tabular baselines.
